% Part: normal-modal-logic
% Chapter: syntax-and-semantics
% Section: relational-models

\documentclass[../../../include/open-logic-section]{subfiles}

\begin{document}

\olfileid[fr]{nml}{syn}{rel}

\olsection{Modèles relationnels}

La notion fondamentale de la sémantique des logiques modales normales
est celle de \emph{modèle relationnel}. Un tel modèle comprend un
ensemble de mondes reliés par une « relation d'accessibilité »
binaire et une valuation qui précise, pour chaque
!!{propositional variable}, dans quels mondes elle est « vraie ».

\begin{defn}
  Un \emph{modèle} pour le langage modal de base est un triplet
  $\mModel{M} = \tuple{W, R, V}$, où
  \begin{enumerate}
  \item $W$ est un ensemble non vide de « mondes »;
  \item $R$ est une relation binaire d'accessibilité sur~$W$;
  \item $V$ est une fonction qui attribue à chaque
    !!{propositional variable}~$p$ un ensemble~$V(p)$ de mondes
    possibles.
  \end{enumerate}
  Lorsque $Rww'$ vaut, nous disons que $w'$ est \emph{accessible
    depuis}~$w$. Lorsque $w \in V(p)$, nous disons que $p$ est
  \emph{vraie en}~$w$.
\end{defn}

La sémantique relationnelle présente l'avantage de permettre des
diagrammes simples, comme celui de la~\olref{fig:simple}. Les mondes y
sont représentés par des nœuds; le monde~$w'$ est accessible
depuis~$w$ si et seulement si une flèche va de~$w$ à~$w'$.
De plus, nous marquons un nœud par~$\mTrue{p}$ si $w \in V(p)$,
et par~\mFalse{p} sinon. La \olref{fig:simple} représente le modèle
où $W = \{w_1, w_2, w_3\}$, $R = \{\tuple{w_1, w_2},\tuple{w_1,
  w_3}\}$, $V(p) = \{w_1, w_2\}$ et $V(q) = \{w_2\}$.

\begin{figure}
  \begin{center}
    \begin{tikzpicture}[modal]
      \node[world] (w1) [label={[align=right]right:\mTrue{p}\\ \mFalse{q}}]{$w_1$}; 
      \node[world] (w2) [label={[align=right]right:\mTrue{p}\\ \mTrue{q}},
        above right=of w1]{$w_2$}; 
      \node[world] (w3) [label={[align=right]right:\mFalse{p}\\ \mFalse{q}},
        below right=of w1] {$w_3$};
      \draw[->] (w1) to (w2);
      \draw[->] (w1) to (w3);
    \end{tikzpicture}
  \end{center}
  \caption{Un modèle simple.}
  \ollabel{fig:simple}
\end{figure}

\end{document}
